% answering.tex -- Module 6: Answering HSC questions (NESA key verbs + model answers).
% Verb meanings are written in our own simple words (not the NESA glossary text).
% Model answers: compare (Ch 1), explain (Ch 3), analyse (Ch 4), assess (Ch 2),
%   F = 0.20 x 3.0 x 0.50 = 0.30 N;  10 MW at 20 kV: I = 500 A, loss 1.25e6 W (12.5 %);
%   at 200 kV: I = 50 A, loss 1.25e4 W (0.125 %);  wires 2.0 A and 6.0 A, 5.0 cm:
%   B1 = 8.0e-6 T, B2 = 2.4e-5 T, F/l = 4.8e-5 N/m on each.
\unnumberedchapter[\faPenNib]{Answering HSC questions}

% marking-point tag used in the model answers below
\newcommand\ansPt[1]{\textcolor{hscbrand}{\bfseries(#1)}}

Every HSC question starts with a \textbf{key verb}, such as ``explain'' or ``evaluate''.
The verb tells you what kind of answer the marker wants (see the table below). Then count the
marks: give about \textbf{one correct point for each mark}. Write one idea in each sentence,
and link the ideas with ``so'' or ``because''. Use physics words, formulas and labelled
diagrams. In Module~6, always give a direction as well as a size. Name the rule you used:
grip rule, palm rule or Lenz's law. In the model answers below, each point that earns a mark
is numbered, like \ansPt{1}.

\begin{table}
  \caption*{The key verbs used in HSC Physics questions}
  \centering
  \small
  \begin{tabular}{@{}L{19mm}L{58mm}L{64mm}@{}}
    \toprule
    \colhead{Verb} & \colhead{What it asks you to do} & \colhead{Module 6 example}\\
    \midrule
    \textbf{Identify} & Name the thing. A word or a short phrase is enough.
      & Identify the part of a DC motor that reverses the current.
        \emph{(The split-ring commutator.)}\\
    \textbf{Outline} & Give the main points only, in order. No details.
      & Outline how a transformer steps up a voltage. \emph{(AC; changing flux in the
        core; emf in the secondary; more secondary turns.)}\\
    \textbf{Describe} & Say what it is like or what happens. You do not need to say why.
      & Describe the path of an electron moving at \ang{90} to a uniform magnetic field.
        \emph{(A circle at a constant speed.)}\\
    \textbf{Explain} & Say \emph{why} or \emph{how} it happens. Link each cause to its
      effect with ``because'' or ``so''.
      & Model answer~2 below.\\
    \textbf{Compare} & Give the similarities \emph{and} the differences.
      & Model answer~1 below.\\
    \textbf{Analyse} & Break it into parts (forces, fields, steps, numbers). Show how they
      are linked and what follows.
      & Model answer~3 below.\\
    \textbf{Assess} & Make a judgement (is it true? how good? how big?) and give reasons.
      & Model answer~4 below.\\
    \textbf{Evaluate} & Judge it against criteria: give strengths \emph{and} limitations,
      then a clear final judgement.
      & Model answer~5 below.\\
    \textbf{Justify} & Give reasons or evidence that support an answer or a decision.
      & Justify the use of laminated transformer cores. \emph{(They break up eddy
        currents, so less energy becomes heat.)}\\
    \textbf{Calculate} & Work out a number. Show the formula, the substitution with units,
      and the answer with a unit.
      & \qty{0.20}{m} of wire, \qty{3.0}{A}, at \ang{90} to \qty{0.50}{T}:
        $F = lIB\sin\theta = \qty{0.30}{N}$.\\
    \bottomrule
  \end{tabular}
\end{table}

\section*{Model answer 1: compare (3 marks)}

\emph{An electron enters a uniform electric field at \ang{90} to the field lines. Another
electron enters a uniform magnetic field at \ang{90} to the field lines. Compare their
motions.}

\begin{modelanswer}
\ansPt{1} \textbf{Similarity:} both fields push on the electron with a force proportional to
its charge ($F = qE$ and $F = qvB$). So in both fields its path curves.

\ansPt{2} \textbf{Difference (path):} the electric force is constant in size and direction
(against the field lines). So the path is a parabola, like a projectile. The magnetic force
is always at \ang{90} to the velocity. So the path is a circle.

\ansPt{3} \textbf{Difference (speed):} the electric field does work on the electron
($W = qEd$), so its speed increases. The magnetic force does no work, so only its direction
changes.
\end{modelanswer}

\begin{markingcriteria}
  \criterion{Points (1), (2) and (3): a similarity and two differences, with reasons}{3}
  \criterion{A similarity and one difference, or two differences}{2}
  \criterion{One correct feature of either motion}{1}
\end{markingcriteria}

\section*{Model answer 2: explain (3 marks)}

\emph{A strong magnet is dropped down a long vertical copper pipe. It takes much longer to
reach the bottom than a non-magnetic object of the same size and mass. Explain why.}

\begin{modelanswer}
\ansPt{1} As the magnet falls, the flux through the copper rings near it changes. By
Faraday's law, this induces eddy currents around the pipe.

\ansPt{2} By Lenz's law, the eddy currents oppose the change in flux. The copper below the
magnet repels it, and the copper above attracts it: an upward force.

\ansPt{3} This force grows with speed until it equals the weight. Then the magnet falls at
a slow, steady speed, and its potential energy becomes heat.
\end{modelanswer}

\begin{markingcriteria}
  \criterion{Points (1), (2) and (3), linked}{3}
  \criterion{Two of these points, linked}{2}
  \criterion{One relevant point, for example ``currents are induced in the copper''}{1}
\end{markingcriteria}

\section*{Model answer 3: analyse (4 marks)}

\emph{A power station sends \qty{10}{MW} through lines with a total resistance of
\qty{5.0}{\ohm}. Analyse how the power lost in the lines changes when the voltage is
raised from \qty{20}{kV} to \qty{200}{kV}.}

\begin{modelanswer}
\ansPt{1} \textbf{Current.} From $P = VI$, $I = P/V$. At \qty{20}{kV}:
$I = \qty{10e6}{W} \div \qty{20e3}{V} = \qty{500}{A}$. At \qty{200}{kV}: $I = \qty{50}{A}$.
Ten times the voltage gives one-tenth of the current.

\ansPt{2} \textbf{Loss.} $P_\mathrm{loss} = I^2R$. At \qty{20}{kV}:
$(\qty{500}{A})^2 \times \qty{5.0}{\ohm} = \qty{1.25e6}{W}$ (\qty{12.5}{\percent} of the
power sent). At \qty{200}{kV}: $(\qty{50}{A})^2 \times \qty{5.0}{\ohm} = \qty{1.25e4}{W}$
(only \qty{0.125}{\percent}).

\ansPt{3} \textbf{Relationship.} Loss $\propto I^2$, so $\frac{1}{10}$ of the current gives
$\frac{1}{100}$ of the loss.

\ansPt{4} \textbf{Consequence.} High-voltage lines waste far less energy. Transformers step
the voltage up at the power station and down near homes (this needs AC).
\end{modelanswer}

\begin{markingcriteria}
  \criterion{Points (1) to (4): both currents, both losses, the $I^2$ link and a consequence}{4}
  \criterion{Correct currents and losses, with little analysis}{3}
  \criterion{A correct current or loss, or ``a higher voltage gives a smaller loss''}{1--2}
\end{markingcriteria}

\section*{Model answer 4: assess (3 marks)}

\emph{Two long parallel wires are \qty{5.0}{cm} apart. Wire~1 carries \qty{2.0}{A} and
wire~2 carries \qty{6.0}{A}. A student says: ``Wire~2 feels the bigger force, because its
current is bigger.'' Assess this claim.}

\begin{modelanswer}
\ansPt{1} The claim is \textbf{incorrect}. Both wires feel forces of the same size.

\ansPt{2} Wire~1 makes $B_1 = \mu_0 I_1/(2\pi r) = \qty{8.0e-6}{T}$ at wire~2, so wire~2
feels $F/l = I_2B_1 = \qty{4.8e-5}{N.m^{-1}}$. At wire~1, $B_2$ is three times as big, but
$I_1$ is one-third as big. So $F/l = I_1B_2 = \qty{4.8e-5}{N.m^{-1}}$ again.

\ansPt{3} This must be true by Newton's third law: the two forces are an action--reaction
pair, equal in size and opposite in direction.
\end{modelanswer}

\begin{markingcriteria}
  \criterion{Points (1), (2) and (3): judgement, equal forces shown, Newton's third law}{3}
  \criterion{Judgement with one correct supporting reason}{2}
  \criterion{One relevant point, for example ``$F/l$ depends on $I_1 I_2$''}{1}
\end{markingcriteria}

\section*{Model answer 5: evaluate (4 marks)}

\emph{Some roller coasters use magnetic brakes: metal fins on each car pass between strong
permanent magnets. Evaluate the use of magnetic brakes to stop the cars.}

\begin{modelanswer}
\ansPt{1} \textbf{How it works:} as a fin moves between the magnets, the flux through it
changes. So eddy currents are induced (Faraday's law). By Lenz's law, the force on these
currents opposes the motion. The car's kinetic energy becomes heat in the fins.

\ansPt{2} \textbf{Strengths:} nothing touches, so nothing wears out, and the brake works in
the wet. Permanent magnets need no power, so the brake still works in a power failure.

\ansPt{3} \textbf{Limitation:} the braking force is large at high speed but falls to zero
as the car stops, because the flux stops changing. So the brake cannot stop the car
completely or hold it still.

\ansPt{4} \textbf{Judgement:} magnetic brakes are an excellent, reliable way to slow a fast
car. But they must be used with friction brakes, which finish the stop and hold the car.
\end{modelanswer}

\begin{markingcriteria}
  \criterion{Points (1) to (4): how it works, strengths, a limitation, a judgement}{4}
  \criterion{Strengths and a limitation with a judgement, but little physics}{3}
  \criterion{A strength or a limitation of magnetic brakes}{1--2}
\end{markingcriteria}
